\section{The planar boundary obstruction}\label{sec:planar}

We now combine the fixed-space theorem with the phase inequality.  The
result is a statement about a planar surface assembled from copies of the
relators: if every boundary component but one reads the identity in $G$,
then the remaining component does too.  This extends the one-relator
boundary argument of \cite[Theorem~4.1 and Section~4]{OpenAIKervaire2026};
the new fixed-space input allows the different relator types to be treated
simultaneously.

Cyclically conjugating the words does not change their normal closure or
their exponent sums.  We may therefore write
\begin{equation}\label{eq:word-slots}
 w_i=x_{\nu_{i,1}}^{s_{i,1}}g_{i,1}\cdots
     x_{\nu_{i,N_i}}^{s_{i,N_i}}g_{i,N_i},
 \qquad
 g_{i,t}\in G,\quad s_{i,t}\in\{1,-1\},\quad
 1\leq \nu_{i,t}\leq n.
\end{equation}
Here each occurrence is a single variable letter, coefficients may be the
identity, and $N_i\geq1$ because the $i$th row of $A$ is nonzero.

A \emph{word disk of type $i$ and sign $+$ or $-$} is an oriented disk
whose boundary reads $w_i$ or $w_i^{-1}$, respectively.  Its variable
letters occupy disjoint closed intervals called \emph{slots}, separated
by nonempty \emph{corner arcs} carrying the coefficients.  On both signs
of disk, index the slots by the positions $t$ in \eqref{eq:word-slots}.
Indices are cyclic modulo $N_i$.  The boundary data are then
\begin{equation}\label{eq:slot-conventions}
\begin{array}{c|c|c|c}
 \text{disk sign}&\text{successor of }t&
 \text{variable at }t&\text{corner after }t\\ \hline
 +&t+1&x_{\nu_{i,t}}^{s_{i,t}}&g_{i,t}\\
 -&t-1&x_{\nu_{i,t}}^{-s_{i,t}}&g_{i,t-1}^{-1}
\end{array}
\end{equation}
The exponent $s_{i,t}$ or $-s_{i,t}$ in this table is the slot's
\emph{traversal sign}.  Attach rectangular bands to pairs of slots so
that the orientations extend across the bands.  Every slot is used once,
and paired slots must have the same variable index and opposite traversal
signs.  The remaining boundary consists of corner arcs and band sides;
its word in $G$ is obtained by reading the corner labels, with no label
contributed by a band side.  Changing the starting corner conjugates this
word, so its triviality is well defined.

\begin{proposition}[Planar boundary obstruction]\label{prop:planar}
Let $w_1,\ldots,w_m\in G*F_n$ have exponent-sum matrix of row rank $m$
over $\mathbb Q$.  Let $\Sigma$ be a connected oriented surface of genus
zero obtained from a nonempty finite collection of signed word disks for
these words by attaching bands that pair every slot exactly once, always
with the same variable index and opposite traversal signs.  If the words
on all boundary components except possibly one are trivial in $G$, then
every boundary word is trivial in $G$.
\end{proposition}

\begin{proof}
The proof estimates one unitary operator in two ways.  Mixing disks of
the same type gives an upper bound for its phase.  Reading the operator
around the boundary gives an exact phase formula.  A choice of scalar
phases at the end makes the two expressions differ only by the phase of
the possibly nontrivial boundary word.

\smallskip\noindent
\emph{Disk counts and boundary cycles.}
Let $p_i$ and $m_i$ be the numbers of positive and negative disks of type
$i$.  Each band pairs opposite occurrences of the same variable, whence
\[
 \sum_i(p_i-m_i)a_{ij}=0\qquad(1\leq j\leq n).
\]
Row independence gives $p_i=m_i$.  In the rest of this proof, let $I$ be
the set of types that occur and write $k_i=p_i=m_i\geq1$ for $i\in I$.
Set
\begin{equation}\label{eq:disk-band-counts}
 K=\sum_{i\in I}k_i,\qquad b=\sum_{i\in I}N_i k_i.
\end{equation}
Thus $\Sigma$ has $2K$ disks, $2b$ slots, and $b$ bands.

Let $\mathcal D$ be the set of slots.  The successor permutation
$\sigma$ follows the oriented boundary of each disk; the fixed-point-free
involution $\theta$ pairs the ends of each band.  For $d\in\mathcal D$,
write $c_d\in G$ for the label on the corner from the end of $d$ to the
start of $\sigma d$, as specified in \eqref{eq:slot-conventions}.
Starting at the start of slot $d$, a boundary path crosses a band side to
the end of $\theta d$, then follows its corner to the start of
$\sigma\theta d$; see Figure~\ref{fig:boundary-step}.  Consequently the
cycles $z$ of $\sigma\theta$ are exactly the boundary components.  If
$z=(d_1,\ldots,d_a)$ is written in this order, put
\begin{equation}\label{eq:boundary-word}
 a(z)=a,\qquad r_z=c_{\theta d_1}\cdots c_{\theta d_a}.
\end{equation}
The nonempty corner arcs ensure that there is at least one boundary
component.  If $q$ denotes their number, genus zero gives
\begin{equation}\label{eq:planar-euler}
 2-q=\chi(\Sigma)=2K-b,\qquad \sum_z a(z)=2b.
\end{equation}

\begin{figure}[tb]
\centering
\begin{tikzpicture}[x=0.93cm,y=0.93cm,>=Stealth,
  every node/.style={font=\small}]
  \fill[black!4] (0,0) rectangle (4,1.15);
  \draw (0,1.15)--(4,1.15);
  \draw[very thick,black!55] (0,0)--(0,1.15);
  \draw[very thick,black!55] (4,0)--(4,1.15);
  \draw[->,thick] (0,0)--(2.2,0);
  \draw[thick] (2.2,0)--(4,0);
  \draw[->,thick] (4,0) .. controls (4.35,-0.8) and
    (5.10,-0.95) .. (5.75,-0.48);
  \draw[very thick,black!55] (5.75,-0.48) .. controls
    (6.0,-0.30) and (6.2,-0.10) .. (6.35,0.15);
  \draw[dashed] (4,1.15)--(4.20,1.65);
  \draw[dashed] (6.35,0.15)--(6.60,0.65);
  \draw[dashed] (0,1.15)--(-0.20,1.65);
  \fill (0,0) circle (1.6pt);
  \fill (4,0) circle (1.6pt);
  \fill (5.75,-0.48) circle (1.6pt);
  \node[left] at (0,0.65) {$d$};
  \node[right] at (4,0.65) {$\theta d$};
  \node at (2,0.85) {band};
  \node[above] at (1.8,0.08) {band side};
  \node[below left] at (0,-0.05) {start of $d$};
  \node[below left] at (4,-0.12) {end of $\theta d$};
  \node[below] at (5.25,-1.02) {corner $c_{\theta d}$};
  \node[right,align=left] at (5.9,-0.65)
    {start of\\$\sigma\theta d$};
\end{tikzpicture}
\caption{One boundary step, with the attached slots drawn as thick gray
intervals.  The arrow crosses a band side and then follows the corner
after the partner slot.  This gives the permutation $\sigma\theta$ and
the label $c_{\theta d}$.  The drawing is local: the disk signs and types
are unspecified, and the two band ends may lie on the same disk.}
\label{fig:boundary-step}
\end{figure}

\smallskip\noindent
\emph{Distributing phases among boundary components.}
For any real numbers $(\eta_z)_z$ with $\sum_z\eta_z=0$, we can choose
real numbers $(\gamma_d)_{d\in\mathcal D}$ such that
\begin{equation}\label{eq:boundary-flow}
 \gamma_{\theta d}=-\gamma_d,\qquad
 \sum_{d\in z}\gamma_d=\eta_z.
\end{equation}
To see this, form a graph whose vertices are the boundary cycles and
whose edges are the pairs $\{d,\theta d\}$, with an end incident to the
cycle containing that slot.  Loops and multiple edges are allowed.
Connectedness of $\Sigma$ means that $\sigma$ and $\theta$ act
transitively on $\mathcal D$.  Since $\sigma\theta$ and $\theta$ generate
the same permutation group, this graph is connected.  Choose a spanning
tree and set the values on all other edges to zero.  At a leaf, set the
value at its end of the remaining edge equal to its prescribed sum and
the value at the other end to its negative; remove the leaf and adjust
the prescribed sum at its neighbor.  The zero total ensures the last
vertex is satisfied.  This proves \eqref{eq:boundary-flow}, including
the one-vertex case.

Use the algebra $M_{2b}(\mathcal M)$ of Section~\ref{sec:phase}, acting on
$\mathbb C^{\mathcal D}\otimes\ell^2(G)$, and define
\begin{equation}\label{eq:surface-operators}
 \begin{split}
 S(e_d\otimes v)&=e_{\sigma d}\otimes L(c_d^{-1})v,\\
 \Theta(e_d\otimes v)&=e^{2\pi i\gamma_d}
                      e_{\theta d}\otimes v.
 \end{split}
\end{equation}
Both operators are unitary, and \eqref{eq:boundary-flow} gives
$\Theta^2=I$.  We identify scalar matrices with their tensor products
with the identity on $\ell^2(G)$; their trace $\tau_{2b}$ is therefore
the ordinary matrix trace.

\smallskip\noindent
\emph{The upper phase bound.}
We will choose a scalar unitary involution $J$ and use the factorization
$S\Theta=(SJ)(J\Theta)$.  We arrange that $SJ$ is an involution with
known phase, and use the fixed-space theorem to give $J\Theta$ many
fixed vectors.
For each $i\in I$, order the $k_i$ disks of each sign.  At each position
$t$, their slots form two coordinate spaces, each identified with
$\mathbb C^{k_i}$.  Given $X_i\in\mathrm U(k_i)$, define a scalar unitary
involution $J$ by sending the positive array at $(i,t)$ to the negative
array by $X_i$, and the negative array back by $X_i^{-1}$.
The same $X_i$ is used at every position of type $i$.

The corner conventions give
\begin{equation}\label{eq:reversal-involution}
 JSJ=S^{-1}.
\end{equation}
Indeed, on positive arrays $S$ moves from $t$ to $t+1$ with coefficient
$L(g_{i,t}^{-1})$, whereas on negative arrays it moves from $t$ to
$t-1$ with coefficient $L(g_{i,t-1})$.  Applying $J$ on either side
interchanges these rules; its factors $X_i$ and $X_i^{-1}$ cancel
because the coefficients are constant across each array and commute
with scalar matrices.  Thus $SJ$ is a unitary involution.  It exchanges
the two disk signs, so its diagonal blocks, and hence its trace, are
zero.  Its $-1$ spectral projection $(I-SJ)/2$ has trace $b$, giving
\begin{equation}\label{eq:disk-involution-phase}
 \ell(SJ)=\frac b2.
\end{equation}

We now choose the $X_i$ so that $J\Theta$ has many fixed vectors.  Regroup
the scalar coordinates by variable index $j$ and traversal sign into
spaces $V_{j,+}$ and $V_{j,-}$.  For each occurrence $(i,t)$ with
$\nu_{i,t}=j$, each space has one block of dimension $k_i$: the positive
or negative disk array is selected according to its traversal sign.
In particular both spaces have dimension
\[
 d_j=\sum_{\substack{i\in I,\ 1\leq t\leq N_i\\\nu_{i,t}=j}}k_i,
 \qquad \sum_j d_j=b.
\]
Both $J$ and $\Theta$ exchange $V_{j,+}$ and $V_{j,-}$.  Under the array
coordinates, let $P_j\in\mathrm U(d_j)$ represent
$\Theta:V_{j,+}\longrightarrow V_{j,-}$.  On the return map
$J:V_{j,-}\longrightarrow V_{j,+}$, the block at $(i,t)$ is
$X_i^{-s_{i,t}}$.  Hence
\begin{equation}\label{eq:fixed-space-application}
 (J\Theta)|_{V_{j,+}}
   =\operatorname{diag}_{(i,t):\,\nu_{i,t}=j}
           (X_i^{-s_{i,t}})P_j.
\end{equation}
The signed block counts here are
$\sum_{t:\nu_{i,t}=j}(-s_{i,t})=-a_{ij}$.  The matrix formed by the rows
of $-A$ indexed by $I$ has full row rank, so
Theorem~\ref{thm:fixed-space} supplies $X_i$ for which
\[
 \dim_{\mathbb C}\ker\bigl((J\Theta-I)|_{\bigoplus_j V_{j,+}}\bigr)
 \geq K.
\]
There are equally many fixed vectors on $\bigoplus_jV_{j,-}$: $J$
interchanges the two sums and satisfies
$J(J\Theta)J=(J\Theta)^{-1}$.  Thus the full fixed-space dimension
$f$ of the scalar matrix $J\Theta$ is at least $2K$.
The same conjugacy makes its eigenvalue multiset invariant under
$\lambda\mapsto\lambda^{-1}$.  Since
$h(\lambda)+h(\lambda^{-1})=1$ for $\lambda\ne1$ and is zero at $1$,
\[
 2\ell(J\Theta)=2b-f\leq2b-2K.
\]
Using $S\Theta=(SJ)(J\Theta)$, \eqref{eq:disk-involution-phase} and
Lemma~\ref{lem:phase-subadditivity} now give
\begin{equation}\label{eq:surface-phase-upper}
 \ell(S\Theta)\leq\frac32 b-K.
\end{equation}
This holds for every phase distribution in \eqref{eq:boundary-flow};
the auxiliary choice of $J$ may depend on that distribution.

\smallskip\noindent
\emph{The boundary formula and the cut limit.}
It remains to read the left side of \eqref{eq:surface-phase-upper} from
the boundary.  On the coordinates of a cycle
$z=(d_1,\ldots,d_a)$ of $\sigma\theta$, the operator $S\Theta$ is a
weighted cyclic shift $T_z$ whose step from $d_t$ to $d_{t+1}$ has weight
$e^{2\pi i\gamma_{d_t}}L(c_{\theta d_t}^{-1})$.  The circuit starting
at $d_1$ therefore has weight
\[
 e^{2\pi i\sum_{d\in z}\gamma_d}
 L(c_{\theta d_a}^{-1}\cdots c_{\theta d_1}^{-1})
   =e^{2\pi i\eta_z}L(r_z^{-1}).
\]
Lemma~\ref{lem:cycle-phase} yields the exact formula
\begin{equation}\label{eq:boundary-phase}
 \ell(T_z)=\frac{a(z)-1}{2}
       +\tau\bigl(h(e^{2\pi i\eta_z}L(r_z^{-1}))\bigr).
\end{equation}

Choose the possibly exceptional component $z_0$, and for $\epsilon>0$
set
\[
 \eta_z=-\epsilon\quad(z\ne z_0),\qquad
 \eta_{z_0}=(q-1)\epsilon.
\]
These phases sum to zero, so \eqref{eq:surface-phase-upper} applies.
For $z\ne z_0$, the assumption $r_z=1$ makes the last term of
\eqref{eq:boundary-phase} equal to $1-\epsilon$ when $0<\epsilon<1$.
For $z_0$, the scalar phase approaches zero from the nonnegative side.
Thus $h(e^{2\pi i(q-1)\epsilon}\lambda)\to h(\lambda)$ at every point
of the unit circle, including $\lambda=1$.  Bounded convergence in the
spectral measure gives
\[
 \tau\bigl(h(e^{2\pi i(q-1)\epsilon}L(r_{z_0}^{-1}))\bigr)
 \longrightarrow \tau\bigl(h(L(r_{z_0}^{-1}))\bigr).
\]
This also covers $q=1$, when the scalar phase is identically zero.
Summing \eqref{eq:boundary-phase} over all cycles and using
\eqref{eq:planar-euler}, we obtain
\begin{align*}
 \lim_{\epsilon\downarrow0}\ell(S\Theta)
 &=\frac{2b-q}{2}+q-1+\tau\bigl(h(L(r_{z_0}^{-1}))\bigr)\\
 &=\frac32 b-K+\tau\bigl(h(L(r_{z_0}^{-1}))\bigr).
\end{align*}
The upper bound \eqref{eq:surface-phase-upper} forces the last trace to
be zero.  Faithfulness of $\tau$ and nonnegativity of $h$ imply
$L(r_{z_0}^{-1})=I$.  Applying this operator to $\delta_1$ gives
$r_{z_0}=1$, as required.
\end{proof}
